\documentclass[12pt]{article}
\parindent 0.0in
\parskip 0.1in
\pagestyle{empty}
\def\R{\mbox{I}\negthinspace \mbox{R}}
\topmargin -.5in
\textheight 8.0in
\textwidth 6.0in
\oddsidemargin 0.2in
\usepackage{amsmath}
\usepackage{hyperref}

\begin{document}

\newcommand{\bm}[1]{\mbox{\boldmath$#1$}}

\begin{center}{\bf Homework Set \#2 -- Due 11:59pm, Friday, April 17.} \\ {\sc MATH 447 Spring 2026}\end{center}
\textbf{Note}: 
\begin{itemize}
\item This homework assignment is mainly based on the Chapter 5 material. 
\item This homework assignment is designed to help you understand how to apply cross-validation and bootstrap.
\item For your convenience, all the data files and the .tex file for this problem are included.
\item Your need to submit two files: A PDF file which contains all the answers and an R script file (with an ``.R'' extension) containing all the code. Any answers presented in the R script file only will not be graded.
\item Your answers must be presented in the same order as the problems given in this assignment, including all graphs and {\tt R} outputs, if applicable (i.e., the graphs and outputs must be in the same order and must not be relegated to the back of your assignment). 
\item Use your judgment to include an appropriate amount of {\tt R} code and output in your answers as necessary. For example, it is not necessary for the instructor to see every single observation of the dataset you used.
\item All the {\tt R} code must be well documented and submitted as a single R script file (which has the ``.R'' extension) to Canvas. 
\item If you need an extension, you need to ask for permission before 5pm on the day the homework is due.
\end{itemize}

The grading scheme is as follows:
\begin{itemize}
\item This homework assignment will be graded out of $30$ points.
\item Each problem will be graded out of $10$ points.
\item If the R code in the .R file is missing or incomplete, at most $5$ points will be deducted.
\end{itemize}

\textbf{Textbook Problems}: Refer to \textit{An Introduction to Statistical Learning with Applications in R}, First Edition by James, Witten, Hastie, and Tibshirani.

\begin{enumerate}
\item  \#5.4.3 and \#5.4.4 (These two problems are combined together as one problem.) For (a), explain the algorithm using schematic diagram(s). For 5.4.4, answer this question using $k$-fold cross-validation and bootstrap. The answer should contain two parts. In the first part, describe how to choose the best model using $k$-fold cross-validation, assuming that you have multiple candidate models. In the second part, by applying the bootstrap idea to the residuals of the model, estimate the standard deviation of one prediction given a set of values for the explanatory variables.\\

\item Let $x_{\mbox{\tiny U},i}$ and $x_{\mbox{\tiny M},i}$ be the average daily social media usage hours ({\tt usage}) and mental health score ({\tt mental}), respectively, for the $i$-th person, where $i = 1,\ldots, 1705$. Also, let $y_i$ be sleep hours per night. 
\begin{itemize}
\item[(a)] Present the scatterplot matrix of ($x_{\mbox{\tiny U},i}$ vs.\ $y_i$), ($x_{\mbox{\tiny M},i}$ vs.\ $y_i$), and ($x_{\mbox{\tiny U},i}$ vs.\ $x_{\mbox{\tiny M},i}$), and describe how these three pairs of variables are related to each other. Also, note irregular patterns, if any.\\
\item[(b)] Suppose that 
\begin{eqnarray*}
y_i = \beta_0 + \beta_1 x_{\mbox{\tiny U},i} + \cdots + \beta_p x_{\mbox{\tiny U},i}^p + \varepsilon_i,
\end{eqnarray*}
where $E[\varepsilon_i] = 0$ and $\mbox{Var}(\varepsilon_i) = \sigma^2$.\\

Using the LOOCV, choose an appropriate order $p$ and justify your choice. You may assume that $p \in \{1, 2, \ldots, 10\}$.\\
\item[(c)] Repeat (b), but use $x_{\mbox{\tiny M},i}$ this time.
\item[(d)] Using the model you chose in (b), predict $y_i$ assuming that $x_{\mbox{\tiny U},i}$ is equal to its sample median. To obtain a prediction, simply modify: \\

\small
{\tt predict(model, newdata=data.frame(usage = 0), data = data\_usage)\\}
\normalsize

appropriately by replacing {\tt model} with an appropriate {\tt glm} model you chose in (b), and {\tt 0} with an appropriate value.\\
\item[(e)] (No need to report in your answer, but appropriate R code needs to be presented in your R script file). Generate $1000$ bootstrap resamples of {\tt data\_usage} by randomly choosing rows with replacement. Use {\tt sample()} with {\tt replace = TRUE} in the input argument.
\item[(f)] For each of the $1000$ bootstrap samples you created in (e), obtain a coefficient estimate corresponding to the linear term $\beta_1$ assuming a polynomial model of order $p = 3$. Then, using {\tt quantile()}, calculate the $2.5$-th and $97.5$-th percentile of $\beta_1$ estimates obtained from $1000$ bootstrap samples. That gives a $95\%$ confidence interval for $\beta_1$.
\end{itemize}

\item  \#5.4.8. However, DO NOT use the {\tt R code} presented in (a). Instead, complete (a)--(f) using the following code.\\

{\tt set.seed(50)\\
x <- rnorm(100, sd=1)\\
y <- 2*x-5*x\string^2+0.01*x\string^3+rnorm(100, sd=10)\\}

For (c), use {\tt set.seed(50, sample.kind="Rejection")}. For (d), use {\tt set.seed(4391, sample.kind="Rejection")}. See HW2code.R to complete this problem.

\end{enumerate}
\end{document}